\clearpage\SetRegion{引言}{teal}{compass}
\phantomsection\label{ch1-preface}\pdfbookmark[0]{引言：开源共创智能未来}{ch1-bm-preface}
\kicker{PREFACE / 从愿景进入研究}
\PageTitle{开源共创智能未来}{让开放成果成为可以接续的共同基础}
当一种技术已经可以被下载，为什么仍有许多组织难以用好它？当一个项目获得广泛关注，谁来承担后续的审阅、修复和维护？当模型跨越地域被反复改造，最初的发布者、下游贡献者与实际使用者之间，怎样形成可持续的合作关系？

这些问题构成了本报告的研究起点。开源合作让知识能够被接续，让创新不必总是从零开始；但开放的意义，不止于成果从一个组织流向另一个组织，还在于更多人能否理解它、改进它，并掌握持续使用它的能力。

人工智能扩大了这种合作的对象。除了源代码，模型权重、训练和评估说明、数据、工具链与运行环境，也进入开放与复用的讨论。它们并不具有相同的生产方式、许可条件与使用门槛，不能仅因为可以下载，就被放入同一种“开放”计量。

\topic{以发展问题组织证据，而不以单一榜单概括世界}
本报告沿\textbf{参与、供给、复用、治理与公共价值}展开。生态地图首先定位参与和资源，主体排行进一步识别网络位置，技术全景解释项目之间的功能分工，模型生态则观察供给、派生与开放条件。不同视角相互补充，最终服务于对行动对象、支持条件和后续成效的讨论。

对政府与国际组织而言，这意味着把“资源增长了多少”与“哪些能力需要支持”放在一起考虑；对企业、研究机构与社区而言，则意味着看见原始发布之外的适配、测试、文档和维护工作。报告不预设各地区应采取相同路径，也不以平台热度直接判断技术质量或社会成效。

\topic{宣言表达共同愿景，研究检验具体进展}
本报告以《全球开源与人工智能国际合作宣言》为价值起点。宣言提出时代呼唤、共同愿景、核心原则与行动共识，后续研究则依据明确的对象、来源和方法，观察其中涉及的发展条件。\textbf{愿景不能代替事实，指标也不能穷尽价值。} 二者的联系，建立在可追问、可核对的证据之上。

下一页说明全书的结构分工；随后以一张结构图梳理宣言，再完整收录现行中文文本。第一章由此进入发展目标与观察方法，为后续的全球观测、主体比较、技术分析与模型研究建立共同语言。

\vfill\ReportNote{宣言依现行中文工作稿收录，保留其原有表述。它是一份价值倡议，与本报告的统计发现、研究推断和行动建议分别阅读；发布署名与机构安排另行核对。}
\clearpage
\SetRegion{全书结构与阅读路径}{teal}{compass}
\phantomsection\label{ch1-architecture}\pdfbookmark[0]{全书结构与阅读路径}{ch1-bm-architecture}
\kicker{REPORT ARCHITECTURE / 六章四视角}
\PageTitle{从发展问题，到证据与行动}{六章各有职责，附录提供复核入口}
\noindent 第一章建立问题与方法；第二至第五章分别观察地域生态、关键主体、技术分工与模型复用；第六章依据这些证据提出有条件的行动建议。
\begin{center}\includegraphics[width=\linewidth]{figures/new/report_architecture_full.pdf}\end{center}
\noindent{\sffamily\bfseries\fontsize{10.5}{15}\selectfont 图导-1\quad 全书结构与研究路径}
\src{依据报告宪法与全书主张地图绘制。连线表示研究与阅读关系，不代表已验证的因果关系。四种证据不合成为单一总榜。\SourceCite{https://github.com/will-ww/2026-GOSD-Report}{13}}
\vspace{2mm}
\noindent 需要把握方向时，可从执行摘要与第一章进入；需要查找地区、主体或技术问题时，可直接进入相应章节，并返回附录核对指标定义和来源。
\clearpage
\SetRegion{宣言结构}{teal}{compass}
\phantomsection\label{ch1-declarationmap}\pdfbookmark[0]{宣言结构：四部分、三层七原则}{ch1-bm-declarationmap}
\kicker{DECLARATION / 结构导读}
\PageTitle{从共同愿景，走向行动共识}{全球开源与人工智能国际合作宣言}
\noindent 宣言以四个部分展开；七项原则分属三个层次，从开放互联和安全可信，延伸到伦理治理、共同创造与社会价值。
\begin{center}\includegraphics[width=\linewidth,height=157mm,keepaspectratio]{figures/new/declaration_structure_2026.pdf}\end{center}
\noindent{\sffamily\bfseries\fontsize{10.5}{15}\selectfont 图导-2\quad 宣言的四部分与三层七原则}
\src{图中标题、原则与引语沿用现行2026年宣言；结构图用于导读，不替代随后收录的全文。}
\clearpage



